\subsection{可标定的工程实现}
\label{sec:engineering}

理想权重 $c_{pq}$ 的两个因子在运行时都不可直接获得：$\gamma_{pq}$ 的后验构造依赖一组缺乏直接估计途径的分布参数，$v_q$ 依赖未知的信号光子数。本节给出两个因子的可标定工程形式，保持权重的乘性结构不变。对于 $\gamma_{pq}$，不显式构造后验，而以核函数直接参数化，
\begin{equation}
  g_{\Delta}(\Delta_{pq})=\exp\!\Big(-\frac{\Delta_{pq}^2}{2\sigma_w^2}\Big),
  \;
  g_{\rho}(\rho_{pq})=\rho_{\mathrm{floor}}+(1-\rho_{\mathrm{floor}})\rho_{pq},
  \label{eq:reliability}
\end{equation}
其中 $g_{\Delta}$ 与几何似然比 $\Lambda_{\Delta}$ 同为高斯形状，随 $|\Delta_{pq}|$ 单调递减，尺度参数 $\sigma_w$ 承担 $\sigma_s$ 的角色；$g_{\rho}$ 为线性函数，保留互相关证据似然比``$\rho_{pq}$ 越大支持越强''的单调走向，$\rho_{\mathrm{floor}}\in[0,1)$ 为低分邻居的保底贡献，避免匹配得分偶然偏低时将仍有先验支持的邻居完全排除。归一化系数会吸收常数因子，代理只需复现理想权重的单调行为而无需保留概率语义，参数由此可在验证集上直接标定。第二个因子则是对齐误差的平均统计响应。

对于 $1/(v_q + \sigma_\epsilon^{2})$，$v_q$取式~\eqref{eq:v_classical}的形式，并除掉直方图自身携带的 $\hat N_q$ 倍幅度，得到：
\begin{equation}
  \frac{1}{\hat N_q\,\big(v_q+\sigma_\epsilon^2\big)}
  =\frac{1}{\sigma_h^2+\sigma_\epsilon^2\,\hat N_q}
  \;\propto\;\frac{1}{\hat N_q+\hat N_0},
  \label{eq:precision_proxy}
\end{equation}
其中
\begin{equation}
  \hat N_0\;\triangleq\;\frac{\sigma_h^2}{\sigma_\epsilon^2},
\end{equation}
而 $\hat N_q$ 为直方图峰显著度，即峰值信号窗口 $\Omega_q$ 内扣除背景估计 $\hat\beta_q$ 后的信号光子数。$\hat N_0$ 是使计数方差 $v_q=\sigma_h^2/\hat N_q$ 与对齐方差 $\sigma_\epsilon^2$ 相等的信号光子数。$\hat N_q\ll\hat N_0$ 时计数方差主导，贡献随信号强度近线性增长；$\hat N_q\gg\hat N_0$ 时对齐方差主导，位置精度不再随光子数提升，贡献趋于饱和、系数反比下降，高信号直方图因而不支配融合结果。$\hat N_0>0$ 同时保证零信号邻居的权重不发散。三个标定参数相互耦合较弱，固定搜索窗口半宽 $K$ 与匹配算子后，在验证集上以深度误差为目标逐一维搜索确定并冻结。最终融合权重与归一化系数为
\begin{equation}
  w_{pq}=\frac{g_{\Delta}(\Delta_{pq})\,g_{\rho}(\rho_{pq})}{\hat N_q+\hat N_0},
  \qquad
  \alpha_{pq}=\frac{w_{pq}}{\sum_{r\in\mathcal{G}_p}w_{pr}},
  \label{eq:weight}
\end{equation}
对对齐后的邻域直方图加权融合
\begin{equation}
  F_p^{\mathrm{fuse}}(b)
  =\sum_{q\in\mathcal{G}_p}\alpha_{pq}\,\widetilde F_{pq}(b),
  \label{eq:fuse}
\end{equation}
再由融合直方图提取回波中心，得到最终深度估计 $\hat\mu_p$。整个流程仅使用单目先验与当前 SPAD 观测即可运行，同时保留了理想贝叶斯模型的核心结构。算法~\ref{alg:pipeline}汇总了完整流程。管线分为两级：场景级预计算（第~1--2~行）每场景仅运行一次，产生单目先验 $b^M$；像素级处理（第~3--17~行）完全局限于 $\mathcal{G}_p\cup\{p\}$ 内部，依次完成峰显著度统计、Coates 校正、先验差分对齐与窗搜索、三因子加权，最后归一化融合并提取峰位。除场景级预计算外，各像素之间无数据依赖，整个管线单遍前向、无迭代，计算量随邻域尺寸线性增长，适合并行实现。

\begin{algorithm}[!t]
\caption{单目引导的直方图融合测距（以像素 $p$ 为例）}
\label{alg:pipeline}
\begin{algorithmic}[1]
\small
\Require 场景内各像素的原始直方图 $\{H_q^{\mathrm{raw}}\}$
\Ensure 像素 $p$ 的深度估计 $\hat\mu_p$
\State $d^{\mathrm{rel}}\leftarrow\mathrm{DAV2}\Big(\mathrm{norm}\big(\ln(1+\sum_b H^{\mathrm{raw}})\big)\Big)$ \Comment{单目相对深度}
\State $b^M\leftarrow\mathrm{Med}_{3\times3}\big(\mathrm{LSQ}(d^{\mathrm{rel}})\big)$ \Comment{最小二乘标定至 bin 域}
\For{$q\in\mathcal{G}_p\cup\{p\}$}
  \State $\hat N_q\leftarrow\sum_{b\in\Omega_q}\big(H_q^{\mathrm{raw}}(b)-\hat\beta_q\big)$ \Comment{原始域峰显著度}
  \State $F_q\leftarrow\mathrm{Coates}(H_q^{\mathrm{raw}};M)$ \Comment{式~\eqref{eq:coates}}
\EndFor
\For{$q\in\mathcal{G}_p\cup\{p\}$} \Comment{$q=p$ 时 $\Delta_{pq}=0$、$\rho_{pq}=1$}
  \State $\Delta_{pq}\leftarrow b_p^M-b_q^M$，\quad $k_{pq}^M\leftarrow\mathrm{round}(\Delta_{pq})$
  \For{$\delta\in\{-K,\dots,K\}$}
    \State $\rho_{pq}(\delta)\leftarrow\big\langle F_p,T_{k_{pq}^M+\delta}F_q\big\rangle\big/(\|F_p\|\,\|F_q\|)$ \Comment{式~\eqref{eq:nmf}}
  \EndFor
  \State $\hat\delta_{pq}\leftarrow\argmax_\delta\rho_{pq}(\delta)$, \quad $\hat k_{pq}\leftarrow k_{pq}^M+\hat\delta_{pq}$
  \State $\widetilde F_{pq}\leftarrow T_{\hat k_{pq}}F_q$, \quad $\rho_{pq}\leftarrow\rho_{pq}(\hat\delta_{pq})$
  \State $w_{pq}\leftarrow g_\Delta(\Delta_{pq})\,g_\rho(\rho_{pq})\big/(\hat N_q+\hat N_0)$ \Comment{式~\eqref{eq:weight}}
\EndFor
\State $\alpha_{pq}\leftarrow w_{pq}\big/\sum_{r\in\mathcal{G}_p\cup\{p\}}w_{pr}$ \Comment{式 \eqref{eq:weight}}
\State $F_p^{\mathrm{fuse}}\leftarrow\sum_q\alpha_{pq}\,\widetilde F_{pq}$ \Comment{式\eqref{eq:fuse}}
\State $\hat\mu_p\leftarrow\operatorname{Peak}\big(\mathrm{MF}_{\sigma_h}(F_p^{\mathrm{fuse}})\big)$ \Comment{匹配滤波后取亚 bin 峰位}
\State \Return $\hat\mu_p$
\end{algorithmic}
\end{algorithm}
